\chapter{Federation Directives Quick Reference}
\label{app:directives}

\section{Read the Version Before the Name}

The entries distinguish the Federation v2.6 vocabulary from later additions.
They describe directive contracts, not blanket support by this book's
Hot Chocolate/Cosmo combination. Compare the subgraph's exported declaration
with the composer's supported definition. Chapter~\ref{ch:composition} shows
two products disagreeing about the same directive name.

The version boundaries follow Apollo's federation changelog
~\autocite{apollofederationversionsrepair}. The baseline definitions were
cross-checked against Apollo's versioned implementation as well as the
reference documentation~\autocite{apollofederation26repair}. Core GraphQL
directives such as \code{@skip}, \code{@include} and \code{@deprecated}
are not federation directives.

\section{The Baseline Vocabulary}

\subsection*{\code{@link}}

Location: schema, repeatable. Arguments: \code{url}, optional \code{as},
\code{for} and \code{import} in the link specification
~\autocite{apollolinkrepair}. Imports select names from the linked vocabulary.
The Hot Chocolate export used here has a narrower declaration: no
\code{as} or \code{for}, and a string-list import. See
section~\ref{sec:ch05-the-directives} before copying a generic link signature.

\subsection*{\code{@key}}

Location: object or interface, repeatable. Arguments: required
\code{fields}, optional \code{resolvable} defaulting to true.
The field set identifies an entity. A reference-only stub uses false; a
subgraph that must accept entity lookups needs a resolvable key. Interface
entity support requires Federation v2.3. See chapter~\ref{ch:first-subgraph}.

\subsection*{\code{@external}}

Location: field or object. Declares fields supplied from elsewhere for this
subgraph's federation contract. The Apollo baseline implementation also
accepts an optional \code{reason}; generic documentation does not consistently
show it. The book uses the argument-free form with \code{@requires}.
It is not permission to provide an arbitrary second authoritative value.

\subsection*{\code{@requires}}

Location: field. Argument: required \code{fields} field set.
The resolver needs these external values in its representation.
Ratings' \code{feedbackUrl} requires the session title. It does not move
filtering or pagination into Ratings; see
chapters~\ref{ch:second-third-subgraph} and~\ref{ch:cross-seam-paging}.

\subsection*{\code{@provides}}

Location: field. Argument: required \code{fields} field set.
Values are available along this particular return path, not every path to
the entity. The provided field must satisfy the external/shareable rules.
Chapter~\ref{ch:second-third-subgraph} records a mismatch between those
published prerequisites and the pinned composer's acceptance; the final
example does not use this optimization.

\subsection*{\code{@shareable}}

Location: object or field; repeatable from Federation v2.2. No arguments.
Permits multiple resolvers for a field, with equivalent answers required.
A declaration does not prove equivalent behavior or fresh approval from
another team. Shared value types need no entity key.
See chapters~\ref{ch:second-third-subgraph} and~\ref{ch:ownership}.

\subsection*{\code{@override}}

Location: field. Argument: required \code{from}, naming the old subgraph.
Deploy the receiving resolver before routing to it, then retire the old
resolver. The optional progressive-override \code{label} was added in v2.7;
the pinned Cosmo composer rejects it even though the Hot Chocolate attribute
exposes it. See section~\ref{sec:ch17-handing-a-field-over}.

\subsection*{\code{@inaccessible}}

Location: type-system elements except schema; no arguments.
Removes an element from the client contract while retaining it for planning.
For client compatibility, compare the client-facing document.
That is \code{graphqlClientSchema} when present, rather than the
planner's \code{graphqlSchema}.
See section~\ref{sec:ch17-writing-it-down-somewhere-else}.

\subsection*{\code{@tag}}

Location: type-system elements, repeatable. Argument: required \code{name}.
Adds metadata; consumers such as contract filters assign its meaning.
Apollo's pinned tag v0.3 implementation includes the schema location, while
the subgraph-reference listing omits it
~\autocite{apollotag26repair}. Do not infer acceptance by a particular
subgraph library from either generic listing.

\subsection*{\code{@authenticated}}

Location: field, object, interface, scalar or enum; no arguments.
Requires an authenticated caller. Federation introduced it in v2.5.
The book measures Cosmo's enforcement after the fetch and separately guards
the field inside Hot Chocolate. It does not secure direct subgraph access
by itself; see chapter~\ref{ch:entities-authorization}.

\subsection*{\code{@requiresScopes}}

Location: the same authorization locations: field, object, interface, scalar
or enum. Argument: required \code{scopes}, a nested list of scope names.
The outer alternatives are OR; the names inside one alternative are AND.
Introduced in v2.5. Runtime configuration and enforcement are not exercised
by the example~\autocite{apolloscopes26repair}.

\subsection*{\code{@policy}}

Location: field, object, interface, scalar or enum. Argument: required
\code{policies}, a nested list of policy names.
Introduced in v2.6. Policy names require an evaluator configured outside the
schema; declaring them is not evidence that such an evaluator ran.
This example does not exercise policy evaluation
~\autocite{apollopolicy26repair}.

\subsection*{\code{@interfaceObject}}

Location: object; no arguments. Introduced in v2.3.
Lets an object declaration stand for an entity interface and contribute
fields to its implementations. It is not an ordinary interface annotation.
The package contains this directive, but the example does not exercise its
execution behavior.

\subsection*{\code{@composeDirective}}

Location: schema, repeatable. Argument: \code{name}, including the
directive's leading \code{@}. Preserves an imported custom directive through
composition; it does not implement that directive's policy.
Apollo's reference signature and pinned compiler differ on the argument's
nullability. Pass an explicit name and verify the chosen tool's declaration.
See chapter~\ref{ch:ownership}.

\subsection*{\code{@extends}}

Location: object or interface; no arguments. An alternative for libraries
that cannot express GraphQL's \code{extend} syntax.
It is not what makes a key resolvable, and this example does not need it.

\section{Later Vocabulary: Reference Only}

These names are outside the example's declared federation version.
Their presence in current documentation does not authorize adding them to
the pinned schema. The following entries are reference contracts
~\autocite{apollodirectives,apollofederationversionsrepair}.

\subsection*{\code{@context}}

Minimum federation version: v2.8. Location: object, interface or union,
repeatable. Argument: required \code{name}. Names an ancestor context.
The versioned implementation confirms repeatability, which one current
reference signature omits~\autocite{apollocontext28repair}.

\subsection*{\code{@fromContext}}

Minimum federation version: v2.8. Location: argument definition.
Argument: \code{field}, a context-field selection. Supplies a resolver
argument from a named ancestor context; it is not a client argument or an
entity key~\autocite{apollocontext28repair}.

\subsection*{\code{@cost}}

Minimum federation version: v2.9. Locations include fields, arguments,
input fields, objects, scalars and enums. Argument: required integer
\code{weight}. The cost annotation emitted by this book's Hot Chocolate
core uses a different definition. Chapter~\ref{ch:composition} diagnoses
that mismatch rather than treating the shared name as interoperability.

\subsection*{\code{@listSize}}

Minimum federation version: v2.9. Location: field.
Its four arguments are:

\begin{minted}{text}
assumedSize: Int
slicingArguments: [String!]
sizedFields: [String!]
requireOneSlicingArgument: Boolean = true
\end{minted}

They describe the estimate, the arguments that bound it, the fields it sizes
and whether exactly one slicing argument must be supplied.
These are cost-analysis hints, not a pagination algorithm.
Chapter~\ref{ch:composition} records an extra Hot Chocolate argument that the
pinned Cosmo definition rejects.

\subsection*{\code{@cacheTag}}

Minimum federation version: v2.12. Location: object or field, repeatable.
Argument: required \code{format}. Declares cache-tag metadata.
It is not used or tested by this example; no cache invalidation behavior
should be inferred for the pinned router.

\section{Related, but a Separate Specification}

\subsection*{\code{@connect}}

A Connector directive for mapping a field to an HTTP source. Connectors
began with a Federation v2.10 requirement, but have their own versioned
specification. They are not the reference-resolver mechanism used here
~\autocite{apolloconnectorsrepair}.

\subsection*{\code{@source}}

A Connector directive that configures a named HTTP source for connector
mappings. It belongs with \code{@connect}, not in the v2.6 baseline.
No connector configuration is included in this book's runnable source
~\autocite{apolloconnectorsrepair}.
